\begin{textsample}{NEG Dim 1 – 2022\_c6 – Score -6.00 – 2022\_c6/enade\_2022\_teo\_00a9aa1edd.txt}  \label{ex:f1_neg_183}
CAPÍTULO 7 ANÁLISE TÉCNICA DA PROVA Este capítulo tem por objetivo apresentar o desempenho dos estudantes concluintes de Teologia no Enade/2022.
Para isso, foram calculadas as estatísticas básicas da prova em seu todo ( Subseção 7.1.1 ), bem como as estatísticas relacionadas aos Componentes de Formação Geral ( Subseção 7.1.2 ) e de Componente Específico da Área ( Susbseção 7.1.3 ).
Dadas as suas características, foram analisadas, em separado, as questões objetivas ( Seção 7.2 ) e as questões discursivas ( Seção 7.3 ).
Tomando -se como base as duas questões discursivas do Componente de Formação Geral, nas Subseções 7.3.1.6 a 7.3.1.8, são apresentados os resultados e os comentários sobre a correção das respostas em relação à Língua Portuguesa.
Para as questões objetivas, foram disponibilizados os índices de Facilidade e de Discriminação Ponto-Bisserial, também em separado, para os Componentes de Formação Geral ( Subseção 7.2.1 ) e de Componente Específico ( Subseção 7.2.2 ).
De \textbf{cada} componente, uma das questões foi escolhida para exemplificar a análise gráfica, relacionando -se as alternativas escolhidas pelos estudantes ( inclusive o gabarito ) com o número de acertos no componente.
No Anexo IV, apresenta -se a íntegra da análise gráfica para todas as questões objetivas.
Para \textbf{cada} uma das questões discursivas, os conteúdos dos tipos mais comuns de respostas dos estudantes são apresentados e comparados com o padrão de resposta esperado ( ver Anexo IX ).
Nas tabelas, constam as seguintes estatísticas das notas⁹ : média do desempenho na prova, erro-padrão da média, desvio-padrão, nota mínima, mediana e nota máxima para \textbf{cada} um de seus componentes.
Tais estatísticas contemplam o total de estudantes concluintes da área de Teologia inscritos e presentes à prova do Enade/2022, tendo em vista agregações, ou por grandes regiões e o país como um todo, ou por categoria administrativa, organização acadêmica da IES e modalidade de oferta dos cursos.
Em relação aos gráficos de distribuição de notas, o intervalo considerado foi de 10 unidades, aberto à esquerda e fechado à direita, com exceção do primeiro intervalo, [ 0 ; 10 ], fechado em ambos os extremos.
Para os gráficos de distribuição das notas das questões discursivas, foram consideradas mais duas categorias : questão em branco¹⁰ e nota zero. ⁹ Uma definição dessas estatísticas pode ser encontrada no Glossário inserido no final desse relatório. ¹⁰ Nesse grupo \textbf{estão} incluídas também as respostas classificadas como nulas ou desconsideradas. 7.1 ESTATÍSTICAS BÁSICAS DA PROVA Esta seção apresenta estatísticas selecionadas e histogramas da nota geral ( 7.1.1 ) e de \textbf{cada} componente : Componente de Formação Geral ( 7.1.2 ) e Componente Específico ( 7.1.3 ).
São também apresentadas estatísticas selecionadas de subpopulações, caracterizadas por grande região, categoria administrativa, organização acadêmica e modalidade de oferta. 7.1.1 ESTATÍSTICAS BÁSICAS GERAIS Na Tabela 7.1, são apresentadas as estatísticas básicas da prova, por grande região, dos estudantes concluintes de Teologia.
Como se pode verificar, a partir dos dados mostrados na Tabela 7.1, a média das notas da prova, como um todo, foi 48,9 ( nas seções seguintes serão analisados o Componente de Formação Geral e o Componente Específico ).
Os estudantes da região Nordeste obtiveram a média mais baixa ( 45 ), e os da região Sudeste alcançaram a média mais alta ( 50,4 ).
As demais médias foram : 46,1, na região Centro-Oeste, 47,7 na região Norte, e 49,2, na região Sul.
O desvio-padrão para o Brasil foi de 14,9, sendo o maior desvio-padrão encontrado na região Nordeste ( 16,1 ) e o menor, na região Centro Oeste ( 14,4 ), indicando uma dispersão um pouco menor das notas desta última região.
A região que obteve a maior nota máxima foi a região Sul ( 90,9 ), ao passo que a região que atingiu a menor nota máxima foi a Norte ( 83,8 ).
A mediana do Brasil, como um todo, foi 48,6, sendo a maior mediana obtida na região Sudeste ( 50,3 ), e a menor obtida na região Nordeste ( 44,2 ).
As regiões Centro-Oeste, sudeste e Sul foram as únicas que obtiveram nota mínima zero.
O comportamento das notas dos estudantes de todo o Brasil pode ser observado no Gráfico 7.1, que apresenta um histograma com a distribuição das mesmas : uma distribuição unimodal, com a moda no intervalo [ 30 ; 40 ].
No gráfico 7.1 é possível observar que o intervalo de nota [ 40 ; 50 ] obtiveram a maior percentagem, chegando a 25 \%.
Os estudantes que deixaram a questão em branco ou a nota atribuída a sua resposta foi zero representam 2,5 \%.
Na Tabela 7.2, são apresentadas informações referentes à média da nota final, desagregadas por categoria administrativa, organização acadêmica e modalidade de oferta dos cursos.
Os dados da Tabela 7.2 mostram que, em relação à categoria administrativa, os estudantes das IES privadas obtiveram média 48,9.
O erro-padrão da média para as IES privadas foi de 0,2, indicando uma dispersão extremamente baixa das notas nessa categoria administrativa.
No tocante à organização acadêmica, os estudantes das faculdades obtiveram a média de 49,8, enquanto os de centros universitários e universidades obtiveram, respectivamente, média de 48,3 e 48,6 A média da modalidade a distância ( 48,2 ) foi inferior à média da modalidade presencial ( 50,5 ). 7.1.2 ESTATÍSTICAS BÁSICAS NO COMPONENTE DE FORMAÇÃO GERAL Na Tabela 7.3, são apresentadas as estatísticas básicas em relação ao componente da prova que avalia a Formação Geral dos estudantes concluintes.
A nota de Formação Geral é obtida levando -se em conta os acertos nas oito questões objetivas e a nota média obtida nas questões discursivas 1 e 2, sendo que a nota de \textbf{cada} questão discursiva foi obtida levando -se em conta 80 \% da nota relativa ao conteúdo e 20 \% da nota relativa ao desempenho linguístico.
Conforme os dados da Tabela 7.3, os estudantes de todo o Brasil obtiveram, no Componente de Formação Geral da prova, desempenho médio de 51,1.
Quanto à variabilidade, o desvio-padrão das notas dos estudantes do Brasil foi 18,0.
A maior média foi obtida na região Sudeste ( 52,6 ) e a menor, na região Nordeste ( 47,9 ).
As demais médias foram : 49,0 na região Centro-Oeste ; 50,9, na região Norte ; e 51,0, na região Sul.
Já o maior desvio-padrão foi obtido na região Nordeste ( 18,8 ), e o menor, na região Norte ( 17,1 ).
Os demais desvios foram : 18,3, na região Centro-Oeste ; 17,7, na região Sudeste ; e 17,8, na região Sul.
A maior nota no Componente de Formação Geral da prova do Enade/2022 foi 98,7, obtida por, pelo menos, um estudante na região Sudeste.
A menor nota máxima foi obtida na região Centro-Oeste ( 88,4 ).
A mediana do Brasil foi 51,8 ; sendo que a menor mediana foi encontrada na região Nordeste ( 47,1 ), e a maior, na região Sudeste ( 53,7 ).
A nota mínima nessa parte foi zero em todas as regiões, exceto na região Norte ( 6,2 ).
No Gráfico 7.2, é possível verificar a avaliação do desempenho dos estudantes do Componente de Formação Geral, a partir do histograma da distribuição das notas correspondentes.
Nesse gráfico, a distribuição é unimodal, com moda em ( 50 ; 60 ].
De acordo com o gráfico 7.2, o intervalo de notas que obteve maior porcentagem foi o de [ 50 ; 60 ], que obteve mais de 20 \%.
Os estudantes que deixaram a questão em branco ou cuja resposta gerou a nota zero [ 0 ; 10 ] representam menos que 2,5 \%.
Na Tabela 7.4, são apresentadas as informações referentes ao desempenho dos concluintes do Componente de Formação Geral, em diferentes agregações : categoria administrativa, organização acadêmica e modalidade de oferta.
A partir dos dados da Tabela 7.4, observa -se que as médias dos tipos de organização acadêmica foram : 50,4, para as universidades ; 50,6, para os centros universitários ; e 52,2, para as faculdades.
As IES privadas obtiveram média de 51,0 ; erro-padrão da média de 0,2 e desvio-padrão de 18,0.
No que se refere à modalidade de oferta, os cursos presenciais ( 53,0 ) apresentaram média superior à média dos cursos a distância ( 50,3 ). 7.1.3 ESTATÍSTICAS BÁSICAS DO COMPONENTE DE CONHECIMENTO ESPECÍFICO Na Tabela 7.5, são apresentadas as estatísticas básicas referentes ao Componente Específico da área de Teologia.
A nota do componente de Conhecimentos Específicos leva em conta as notas da parte objetiva e a média das notas obtidas nas três questões discursivas da parte discursiva da área.
De acordo com a Tabela 7.5, a média do desempenho dos estudantes do Brasil, para o Componente Específico da prova, foi 48,1.
A maior média foi obtida na região Sudeste ( 49,6 ), e a menor, na região Nordeste ( 44,1 ).
As demais médias foram : 45,1 na região Centro-Oeste, 46,7 no Norte e 48,6 na região Sul.
Quanto à variabilidade das notas, o desvio-padrão do Brasil foi 15,9 ; sendo o maior desvio-padrão observado nas regiões Norte ( 17,3 ) e Nordeste ( 17,3 ) e o menor, na região Centro Oeste ( 15,2 ).
Os demais desvios foram : 15,1 na região Sul e 16,3 na região Sudeste.
A mediana das notas dos estudantes de todo o Brasil foi 47,7.
A maior mediana ocorreu na região Sudeste ( 49,5 ), e a menor, na região Nordeste ( 43,2 ).
As demais medianas foram : 44,3 na região Centro-Oeste e 44,6 na região Norte e 48,5 Sul.
A nota máxima do Brasil, como um todo, foi 91,3 ; sendo obtida por pelo menos um estudante na região Centro Oeste.
As demais notas máximas foram : 90,4 na região Nordeste ; 90,8 na região Sudeste ; 88,7 na região Norte e 90,1 na região Sul.
A maior nota mínima foi de 14,2, na região Norte.
Assim como nos Gráficos 7.1 e 7.2, mostrados anteriormente, no Gráfico 7.3, apresentado adiante, é possível ser feita uma avaliação do desempenho de concluintes de Teologia, em relação ao Componente Específico, com um histograma da distribuição das notas correspondentes.
Esta também é uma distribuição unimodal, e o grupo modal é o ( 40 ; 50 ], o mesmo do grupo modal para a Formação Geral e para a prova, de forma geral.
O gráfico 7.3 evidencia que o intervalo com maior percentagem é o ( 40 ; 50 ].
Os estudantes que deixaram a questão em branco ou a nota atribuída a sua resposta foi zero [ 0 ; 10 ] representam 1 \%.
Na Tabela 7.6, é apresentada uma comparação dos resultados em relação à categoria administrativa, à organização acadêmica e à modalidade de oferta, levando -se em conta o desempenho de estudantes do Componente Específico da prova.
De acordo com os dados expostos na Tabela 7.6, no que se refere à organização acadêmica, a maior média foi a das faculdades ( 49,0 ), vindo a seguir a dos Centros Universitários ( 47,5 ), seguida pela das Universidades ( 47,9 ).
O maior desvio-padrão foi o das faculdades ( 16,8 ).
As faculdades obtiveram a maior nota máxima ( 91,3 ).
Os centros universitários obtiveram nota máxima 90,3 ; as universidades, 90,4.
As medianas foram : 47,6 nas Universidades, 46,8 nos centros universitários e 49,0 nas faculdades.
Todas as IES obtiveram nota mínima 0.
Quanto à categoria administrativa, as IES privadas abstiveram 48,1 de média ; 15,9 de desvio-padrão e 91,3 de nota máxima. 7.2 ANÁLISE DAS QUESTÕES OBJETIVAS Esta seção apresenta estatísticas selecionadas e histogramas das questões objetivas de Formação Geral ( 7.2.1 ) e do Componente Específico ( 7.2.2 ).
São, também, apresentadas e comparadas as médias das subpopulações caracterizadas por grande região.11 7.2.1 COMPONENTE DE FORMAÇÃO GERAL Na Tabela 7.7, são apresentadas as estatísticas básicas relativas às oito questões objetivas do Componente da prova que abrange a Formação Geral dos estudantes. 11 As fórmulas utilizadas para o coeficiente de correlação ponto-bisserial e para o índice de facilidade \textbf{estão} apresentadas no Capítulo 2.
De acordo com os dados da Tabela 7.7, em relação ao Componente de Formação Geral, a média do Brasil foi 51,2.
A menor média foi encontrada na região Nordeste ( 48,4 ), e a maior, na região Sudeste ( 52,5 ).
As demais médias foram : 49,1 na região Centro-Oeste, 49,3 na região Norte e 51,5 na região Sul.
O desvio-padrão do Brasil foi 21,5 ; sendo o maior desvio-padrão encontrado na região Nordeste ( 22,7 ), e o menor desvio-padrão encontrado na região Centro-Oeste ( 21 ).
As medianas do Brasil, como um todo, e das cinco Grande Regiões foram iguais a 50,0.
A nota máxima ( 100,0 ) foi alcançada nas cinco grandes regiões.
As notas mínimas ( 0,0 ) foram iguais para todas as regiões.
Na Tabela 7.8, exposta mais adiante, são apresentados o Índice de Facilidade e o Índice de Discriminação ( Ponto-Bisserial ) para \textbf{cada} uma das questões objetivas do Componente de Formação Geral.
De acordo com os dados da Tabela 7.8, as questões objetivas do Componente de Formação Geral, segundo o Índice de Facilidade, foram assim avaliadas : nenhuma teve o Índice de Facilidade classificado como Muito fácil.
Quatro questões foram classificadas como fáceis, com valor de 0,61 ; 0,81 ; 0,64 e 0,63.
Duas questões foram classificadas como Médio por terem valores de 0,47 e 0,42 e duas como difíceis com valores 0,25 e 0,28.
Quanto ao índice de discriminação, sete questões foram classificadas como Muito bom, com valores entre 0,40 e 0,52.
Uma questão foi classificada como Bom, com índice de acerto de 0,39.
De modo geral, o Índice de Facilidade variou de 0,25 a 0,81 e o de Discriminação, de 0,39 a 0,52.
Na Tabela 7.9, é apresentada a distribuição das questões, levando -se em conta, simultaneamente, a classificação dos dois índices.
Os dados da Tabela 7.9 atestam que sete questões que tiveram Índice de Discriminação Muito bom figuraram em três níveis de dificuldade – Fácil, Médio e Difícil : quatro classificadas na categoria Fácil, duas classificadas na categoria Médio e uma, na categoria Difícil.
As questões que tiveram Índice de Discriminação Bom figuraram no nível de dificuldade Difícil.
A título de exemplo das análises do comportamento das questões objetivas, no Gráfico 7.4, apresenta -se a análise gráfica da Questão 6 de Formação Geral.
No Gráfico 7.4, \textbf{cada} uma das cinco curvas representa o percentual de respostas em determinada alternativa da questão, em função do número de acertos dos estudantes nessa parte da prova ( Formação Geral/Múltipla Escolha ), antes de possíveis eliminações pelo critério Ponto-Bisserial.
Os gráficos relativos às demais questões de Formação Geral constam do Anexo IV. 7.2.2 COMPONENTE DE CONHECIMENTO ESPECÍFICO Na Tabela 7.10, são apresentadas as estatísticas básicas em relação às questões objetivas do Componente Específico da prova de Teologia por grande região.
Como se verifica a partir dos dados da Tabela 7.10, a média do Brasil para o Componente Específico da Prova foi 50,5.
A menor média foi observada na região Nordeste ( 46,7 ), e a maior, na região Sudeste ( 51,9 ).
O desvio-padrão de todo o Brasil foi 17,0, sendo o maior desvio-padrão encontrado na região Nordeste ( 18,8 ), e o menor, na região Centro Oeste ( 16,2 ).
A mediana de todo o Brasil foi 50,0, o mesmo valor da mediana encontrada nas regiões Sul e Sudeste.
Nas regiões Centro-Oeste, Nordeste e Norte, a mediana foi 45,8.
A nota máxima da prova foi 95,8, obtida nas questões objetivas do Componente Específico, por, pelo menos, um estudante da região Centro-Oeste.
A nota mínima foi zero, obtida por pelo menos um estudante nas regiões Centro-Oeste, Nordeste, Sudeste e Sul.
Na Tabela 7.11, são apresentados os Índices de Facilidade e de Discriminação ( Ponto-Bisserial ) das questões objetivas do Componente Específico, para os estudantes de Teologia.
Para facilitar a diferenciação das questões, serão usadas as mesmas cores da Tabela 7.8, anteriormente apresentada, para as diferentes classificações dos Índices de Facilidade e de Discriminação.
A partir do Índice de Facilidade obtido, apresentado na Tabela 7.11, pode -se concluir que, em relação ao índice de facilidade, a maioria das questões objetivas da prova foram consideradas de nível Médio ( 9 ) e Difícil ( 9 ), oito questões foram consideradas fáceis e uma questão muito difícil.
Quanto ao índice de discriminação, doze questões obtiveram índice bom, oito questões obtiveram índice muito bom, quatro índices médios, e três índices fracos.
De modo geral, o Índice de Facilidade variou de 0,13 a 0,79 e o de Discriminação, de 0,17 a 0,47.
Na Tabela 7.12, é apresentada a distribuição das questões levando -se em conta, simultaneamente, a classificação dos dois índices.
Os dados da Tabela 7.12 atestam que doze questões que tiveram Índice de Discriminação Bom figuraram em três níveis de dificuldade – Fácil, Médio e Difícil : três classificadas na categoria Fácil, cinco classificadas na categoria Médio e quatro, na categoria Difícil.
A título de exemplo das análises do comportamento das questões objetivas, no Gráfico 7.5, analisa -se a Questão 20 do Componente Específico.
No Gráfico 7.5, \textbf{cada} uma das cinco curvas representa o percentual de respostas em determinada alternativa da questão, em função do número de acertos dos estudantes nessa parte da prova ( Componente Específico ), antes de possíveis eliminações pelo critério Ponto-Bisserial.
Os gráficos relativos às demais questões do Componente Específico constam no Anexo IV. 7.3 ANÁLISE DAS QUESTÕES DISCURSIVAS Esta seção apresenta estatísticas selecionadas e histogramas das Questões Discursivas de Formação Geral ( 7.3.1 ) e de Componente Específico ( 7.3.2 ).
Expõe e compara, também, as médias de subpopulações, caracterizadas por grande região.
Cumpre notar que nem todas as questões passam pelo crivo dos corretores.
As respostas depois de digitalizadas passam por um pré-processamento de identificação automático.
Algumas destas questões são definidas como um “ branco automático ” ( que não são enviadas para os professores para correção, a não ser que na segunda fase, a de inspeção visual, seja constatada uma resposta ).
O critério usado neste pré-processamento, é verificar para \textbf{cada} questão, a partir da quantidade de PIXELS, o que provavelmente, é uma resposta em branco.
A partir da separação destes documentos que foram considerados em branco, eles são enviados para uma verificação visual, por uma equipe de colaboradores, devidamente capacitados.
Deste modo, quaisquer erros, deste pré-processamento, são encontrados, e as questões preenchidas são enviadas para a correção dos professores.
Eventualmente, algumas respostas em branco não são detectadas nesse pré-processamento automático.
A presença de pequenos riscos, ou manchas, podem fazer o programa não identificar uma questão como um possível branco.
Essas são enviadas para correção, mas os professores as classificam como “ branco ”. 7.3.1 COMPONENTE DE FORMAÇÃO GERAL As análises dos resultados de desempenho dos estudantes de Teologia, nas duas questões discursivas relativas à Formação Geral, encontram -se na Tabela 7.13 e no Gráfico 7.6, cujos dados são analisados de forma sequenciada.
No caso da Tabela 7.13, analisa -se a média das notas de \textbf{cada} Questão Discursiva de Formação Geral que leva em conta as correções realizadas por duas bancas : uma que analisa e pontua aspectos relativos ao conteúdo das respostas, seguindo o padrão de respostas divulgado pelo Inep ( ver Anexo IX ) ; outra que avalia o desempenho linguístico expresso pela redação das respostas, seguindo um padrão de respostas específico da área de Língua Portuguesa, detalhado na seção 7.3.1.7 deste relatório.
Na Tabela 7.13, observa -se que os estudantes de todo o Brasil obtiveram, nas questões discursivas de Formação Geral média igual a 50,7.
Pode -se notar, também, que o desvio-padrão nesse conjunto de questões foi 23,6.
A maior média foi obtida na região Norte ( 53,3 ) e a menor, na região Nordeste ( 47,2 ).
A mediana de todo o Brasil, neste componente, foi 55,5.
Quanto as regiões, a maior mediana foi a da região Norte ( 60,5 ).
A nota máxima ( 100 ), que foi obtida por pelo menos um estudante na região Sul.
A nota mínima ( 0,0 ) foi a mesma em todas as regiões do Brasil.
No Gráfico 7.6, \textbf{está} representada a distribuição das notas nas questões discursivas do Componente de Formação Geral.
Como mostra o Gráfico 7.6, a moda da distribuição ocorre no intervalo [ 60 ; 70 ], com frequência de aproximadamente 20 \% dos participantes.
No intervalo [ 0 ; 10 ], inclui -se, além da nota zero, a frequência de estudantes que deixaram esse tipo de questão em branco, tem -se um percentual de aproximadamente 10 \%.
Na sequência, os resultados verificados para \textbf{cada} uma das questões discursivas de Formação Geral serão apresentados, estabelecendo -se relações com os conteúdos abordados em \textbf{cada} uma delas.
Os comentários das Bancas de docentes corretores, a respeito do observado na correção das respostas dos estudantes, suas impressões e conclusões serão apresentados junto à análise de \textbf{cada} questão.
Cumpre esclarecer que, tendo em vista que as questões discursivas de Formação Geral são padronizadas, ou seja, constam de todas as provas, os comentários da Banca são os mesmos para todas as carreiras acadêmicas, sendo direcionados a todos os estudantes que participaram do Enade/2022.
A seguir, serão analisados o desempenho linguístico dos estudantes da área de Teologia nas duas questões discursivas de Formação Geral do Enade/2022, apresentando -se os resultados obtidos e os comentários da banca de correção para \textbf{cada} questão. 7.3.1.1 ANÁLISE DE CONTEÚDO DA QUESTÃO DISCURSIVA 1 DO COMPONENTE DE FORMAÇÃO GERAL Os dados de Teologia obtidos a partir das respostas à Questão 1 do Componente de Formação Geral, encontram -se na Tabela 7.14 e no Gráfico 7.7, exposto após a apresentação dos dados da referida tabela.
Verifica -se, a partir dos dados da Tabela 7.14, que na Questão 1 do Componente de Formação Geral – de desempenho abaixo ao da outra questão do mesmo componente –, os estudantes do Brasil todo tiveram média 49,5.
A maior média para a Questão 1 foi obtida na região Norte ( 52,6 ), e a menor, na região Nordeste ( 47,0 ).
Quanto à variabilidade das notas, o desvio-padrão de todo o Brasil foi 26,2.
O menor desvio-padrão foi obtido na região Sudeste e Norte ( 25,5 ), e o maior desvio-padrão foi obtido pela região Nordeste ( 27,2 ).
A mediana do Brasil, como um todo, foi 53,0, sendo igual no Centro-Oeste e Sul.
A nota máxima da Questão Discursiva 1 foi 100,0, obtida por, pelo menos, um estudante de todas as cinco Grande Regiões.
Já a nota mínima para todas as regiões foi zero.
No Gráfico 7.7, apresenta -se a distribuição das notas na Questão Discursiva 1 do Componente de Formação Geral.
O Gráfico 7.7 mostra que mais de 12 \% dos participantes deixaram em branco a resposta à questão discursiva 1.
A moda dessa distribuição ocorre no intervalo ( 50 ; 60 ] com uma frequência maior que 17 \% dos participantes. 7.3.1.2 COMENTÁRIOS SOBRE A CORREÇÃO DE CONTEÚDO DAS RESPOSTAS À QUESTÃO DISCURSIVA 1 A Questão 1 abordou um tema atual, com dois itens para os estudantes responderem : a ) “ discuta por que os investimentos em educação, ciência, tecnologia e inovação são estratégicos e contribuem para o desenvolvimento de um país ” ; b ) “ explique como o fomento público ao desenvolvimento científico pode atender à justiça social e à equidade, em contextos como o da pandemia de Covid-19 ”.
O tema contempla como os investimentos em educação, ciência, tecnologia e inovação ( ECTI ) cooperam para um país se desenvolver e como o fomento público, vinculado ao desenvolvimento científico, pode atender à justiça social, à equidade e à população em momentos de pandemia.
Esse tema é atual, pois o mundo vivenciou recentemente uma pandemia que, somente no Brasil, deixou mais de 700 mil mortos, sendo exaustivamente comentada na televisão e nos meios digitais.
Logo, almeja -se que os respondentes justifiquem a necessidade de investimentos em ECTI com foco em questões como : investimentos público em P\&D por meio de universidades públicas ; pesquisas realizadas em outros países que possibilitaram o acesso mais célere às vacinas e serviram de modelo para o que foi desenvolvido no Brasil ; pesquisadores de universidade, tanto brasileiras quanto estrangeiras como protagonistas para orientar e informar a respeito da Covid-19 por meio dos canais de comunicação ; desenvolvimento de equipamentos de segurança, como máscara de nariz ou de rosto, respiradores, que foram implementados e viabilizados por pesquisa e extensão em universidades e escolas técnicas.
Esperava -se que os candidatos soubessem o básico a respeito do tema, pois o assunto é bastante atual.
Quanto ao item B, esperava -se que os candidatos abordassem questões como : o papel das universidades e dos centros de pesquisa para a justiça e a equidade ; o investimento público em pesquisa, tecnologia e educação para garantir acesso igualitário da população aos recursos de saúde, podendo fazer um contraponto com os investimentos privados ; a política de distribuição de vacinas gratuitas por meio do Sistema Único de Saúde, que possibilitou o acesso às vacinas no Brasil ; e as questões econômicas mundiais que determinaram uma distribuição desigual de recursos de saúde em escala global, com os países mais ricos retendo a maior parte desses recursos.
Atualmente, é comprovado que os países dominantes são aqueles que se desenvolvem e dominam as ferramentas tecnológicas.
As pesquisas mostram que os países que, mesmo em crise, investiram em ciência, tecnologia e inovação são os que conseguiram se reerguer mais rapidamente e sobressair -se aos demais países em crises pandêmicas.
Portanto, a questão foi elaborada de maneira oportuna e adequada ao universo do estudante universitário brasileiro.
No item A, o estudante poderia focar em diversos argumentos, entre eles : como ocorreu o investimento público em pesquisa e desenvolvimento por meio de universidades públicas e institutos de pesquisa ; como as pesquisas realizadas em outros países possibilitaram acesso mais rápido às vacinas e serviram de modelo para o que foi desenvolvido no Brasil ; como o contingenciamento do Fundo Nacional de Desenvolvimento Científico e Tecnológico ( FNDCT ) prejudicou e ainda prejudica o Brasil, entre diversos outros assuntos.
No item B, era esperado que o estudante explicasse como o fomento público ao desenvolvimento científico pode atender à justiça social e à equidade, em contextos como o da pandemia de Covid-19.
Além disso, deveria abordar a importância das universidades com a justiça social ; como os países que investem nessa área têm grandes vantagens em relação a outros países que não investem ; a distribuição de vacina atende à justiça social por possibilitar que toda a população seja vacinada, independentemente de ser rico ou pobre, por exemplo.
O espaço disponível para a resposta dos dois itens foi adequado e a maioria das respostas ocupou parte das linhas, não utilizando todo o espaço.
Algumas respostas não apresentaram a separação dos itens A e B.
A maioria dos estudantes respondeu de modo semelhante ao que constava do padrão de resposta.
Foi recorrente a exposição do problema sem desenvolvê -lo com clareza, apenas de forma superficial.
A abordagem do candidato sobre o tema tornou -se, em alguns casos, muito genérica.
As respostas consideradas boas atenderam ao que a banca esperava.
Abordaram a importância de investir em CT\&I para desenvolver o país, desdobrando os argumentos e citando como isso poderia ocorrer.
No geral, as notas \textbf{estavam} abaixo do esperado.
Ficou evidente que a maioria dos estudantes possui um conhecimento sobre os dois itens apresentados.
A Questão 1 visou avaliar o desenvolvimento da argumentação de um tema relevante para a sociedade brasileira vigente.
Além disso, houve problemas estruturais nos textos, pois a grande maioria dos estudantes tem uma noção do assunto, mas muitos possuem dificuldade de construir um texto com coesão e coerência, ou seja, não conseguem concatenar os argumentos para defender um ponto de vista.
No geral, em relação à Questão 1, identificou -se que a principal deficiência na formação dos estudantes \textbf{está} relacionada ao aprofundamento do conteúdo e à capacidade de expressar, com clareza e coesão, sua opinião a respeito de determinado assunto.
Relatório Síntese de Área – ENADE/2022 7.3.1.3 ANÁLISE DE CONTEÚDO DA QUESTÃO DISCURSIVA 2 DO COMPONENTE DE FORMAÇÃO GERAL A Tabela 7.15 apresenta dados estatísticos a respeito das notas de conteúdo da Questão Discursiva 2 do Componente de Formação Geral, organizados por grande região de funcionamento dos cursos, no país.
A Tabela 7.15 mostra que o desempenho médio dos estudantes, na Questão Discursiva 2 ( média 51,9 ), foi acima do obtido na Questão Discursiva 1 ( média 49,5 ).
A região Nordeste foi aquela cuja média foi menor ( 47,4 ), e a maior média foi na região Sudeste ( 54,2 ).
Quanto à variabilidade das notas, o desvio-padrão de todo o Brasil foi 27,5 ; superior ao obtido na Questão Discursiva 1 ( 26,2 ).
O maior desvio nessa questão foi obtido na região Norte ( 29,9 ), enquanto o menor foi obtido na região Sudeste ( 26,7 ).
A mediana de todo o Brasil foi 57,5.
A nota máxima da Questão Discursiva 2 foi 100,0, obtida por pelo menos em todas as regiões.
Já a nota mínima para todas as regiões foi zero.
No Gráfico 7.8, mostra -se a distribuição das notas na Questão Discursiva 2 do Componente de Formação Geral.
O Gráfico 7.8 mostra que quase 14 \% dos participantes deixaram em branco a resposta à questão discursiva 2.
A moda dessa distribuição ocorre no intervalo ( 70 ; 80 ] com uma frequência maior que 17 \% dos participantes. 7.3.1.4 COMENTÁRIOS SOBRE A CORREÇÃO DE CONTEÚDO DAS RESPOSTAS À QUESTÃO DISCURSIVA 2 A Questão Discursiva 2 abordou um tema importante, com o objetivo de analisar conhecimentos gerais do aluno.
O tema contemplou o conhecimento geral do patriarcado estrutural ( dominação masculina ), o qual é recorrente nos noticiários, e um infortúnio na vida dos brasileiros, pois trata -se de uma questão de violência contra as mulheres.
É recorrente assistir ou ler notícias sobre o domínio masculino na sociedade brasileira.
Questões como a participação da mulher na política ; a paridade de salários entre profissionais de sexos diferentes ; a desigualdade social e o tratamento hostil que é dado às mulheres, em diversas situações, tanto profissionais quanto familiares, devem ser de conhecimento geral do estudante.
Não se objetiva tratar questões de \textbf{ordem} específica para a solução deste problema, mas, apenas, discorrer e argumentar sobre ele.
Portanto, a questão foi elaborada de maneira oportuna e adequada para o universo do estudante universitário brasileiro.
A Questão \textbf{está} dividida em dois itens : A e B.
No item A, esperava -se que o estudante abordasse o patriarcado – dominação masculina que gera a divisão sexual do trabalho, a desigualdade econômica e política, entre homens e mulheres, a sub-representação das mulheres na política e em cargos de poder.
Esperava -se que o estudante ressaltasse, ainda, questões culturais, como a educação sexista e a mídia, que coisificam a imagem das mulheres e naturalizam a desigualdade entre os sexos, de forma a considerar os homens como dominantes, enquanto as mulheres são historicamente tidas como submissas e incapazes.
No item B, era esperado que o estudante apresentasse a importância do papel do Estado na implementação de políticas públicas voltadas à proteção da mulher, tais como na área da saúde, na educação não sexista, na assistência social, na segurança pública – delegacias femininas, defensoria pública especializada, casa de abrigo, juizado especial etc., bem como ONGs que atuam no combate à violência contra a mulher e ao feminicídio.
O espaço disponível para a resposta dos dois itens foi adequado e grande parte das respostas ocupou parte das linhas, não contemplando o espaço completo.
Com relação à linguagem utilizada pelos estudantes, as respostas \textbf{estavam} dentro do esperado.
Alguns estudantes utilizavam termos do jargão do cotidiano.
No item B, o estudante deveria propor duas ações do Estado para o enfrentamento da violência contra as mulheres.
Os alunos deveriam apresentar propostas e desenvolvê -las.
Muitos estudantes responderam apenas por meio da apresentação de ações do Estado, mas sem desenvolvê -las ou Em geral, os alunos apenas expuseram ações – algumas vezes, uma, duas e, até mesmo, três ações –, mas sem fazer relações entre elas ou sem desenvolvê -las.
As respostas dos estudantes demonstraram que a maioria apresenta um conhecimento geral a respeito do patriarcado estrutural e da trágica desigualdade que esta impõe dos homens sobre as mulheres, bem como de possíveis ações do Estado para a proteção delas.
A Questão visou avaliar o desenvolvimento da argumentação de um tema relevante para a sociedade brasileira vigente. 7.3.1.5 CONSIDERAÇÕES FINAIS SOBRE A CORREÇÃO DE CONTEÚDO DAS QUESTÕES DISCURSIVAS DO COMPONENTE DE FORMAÇÃO GERAL As questões discursivas da área de Formação Geral abordaram temas atuais e proporcionaram a avaliação de competências importantes dentre as previstas pelas suas Diretrizes Curriculares.
A análise das respostas revelou, principalmente, dificuldade dos estudantes em relação à seleção de argumentos sólidos para discorrer sobre o tema – que pode ser fruto de desinteresse ou de desconhecimento mais profundo do contexto.
Esses fatos parecem ter sido os principais responsáveis pelo resultado aquém do esperado em termos de avaliação quantitativa.
Mesmo os estudantes que responderam em conformidade com o padrão demonstraram um universo de argumentação muito limitado, pobre em vocabulário, com o recurso frequente de procurar nos textos introdutórios os elementos para a construção de suas frases.
Encadear e relacionar argumentos foi uma deficiência percebida.
As razões eram apresentadas de forma confusa, sem fluência e sem constituir unidade, utilizando termos genéricos e vagos. 7.3.1.6 ANÁLISE DE LÍNGUA PORTUGUESA DAS QUESTÕES DISCURSIVAS DO COMPONENTE DE FORMAÇÃO GERAL Os dados de Teologia obtidos a partir das respostas às questões discursivas do Componente de Formação Geral, no que tange à Língua Portuguesa, são apresentados na Tabela 7.16.
Os dados da Tabela 7.16 mostram que, em relação ao quesito Língua Portuguesa, os estudantes, de todo o Brasil, obtiveram média 59,1.
A maior média em relação à Língua Portuguesa foi obtida na região Norte ( 63,0 ), e a menor, na região Nordeste ( 56,9 ).
Quanto à variabilidade das notas, o desvio-padrão de todo o Brasil foi 26,7.
O menor desvio-padrão foi obtido na região Norte ( 25,6 ), e o maior desvio-padrão foi obtido na região Nordeste ( 29,1 ).
A mediana das notas de Língua Portuguesa foi 65,0 para o Brasil, como um todo, a mesma obtida em todas as regiões.
A nota máxima para todo o Brasil foi 100, com, pelo menos, um estudante tirando essa nota em todas as regiões.
A nota mínima foi zero em todas as regiões do país.
No Gráfico 7.9, mostra -se a distribuição das notas de Língua Portuguesa do Componente de Formação Geral.
O Gráfico 7.9 mostra que mais de 10 \% dos participantes não pontuaram em língua portuguesa nas questões discursivas do componente de formação geral.
A moda dessa distribuição ocorre no intervalo ( 60 ; 70 ] com uma frequência maior que 25 \% dos participantes. 7.3.1.7 COMENTÁRIOS SOBRE A CORREÇÃO DE LÍNGUA PORTUGUESA DAS RESPOSTAS ÀS QUESTÕES DISCURSIVAS Os comentários feitos neste item a respeito da correção de Língua Portuguesa das respostas emitidas para as questões discursivas do Componente de Formação Geral \textbf{estão} agrupados em três eixos : avaliação de desempenho linguístico ; padrão de resposta utilizado na avaliação e análise das respostas dadas pelos estudantes. a ) O padrão de resposta utilizado na avaliação Em relação a este eixo de análise, foram considerados, no padrão de resposta, as seguintes competências : ( 1 ) estruturação textual ; ( 2 ) respeito às convenções ortográficas da norma-padrão da Língua Portuguesa ; ( 3 ) domínio dos diferentes aspectos morfossintáticos próprios da modalidade escrita formal da norma-padrão da Língua Portuguesa ; e ( 4 ) seleção vocabular adequada à modalidade escrita formal da Língua Portuguesa, exigida pela situação comunicativa.
A seguir, são apresentados os aspectos avaliados associados a \textbf{cada} uma destas competências. ( 1 ) Estruturação textual condizente com o gênero textual solicitado e o modo de organização textual adequado ao gênero ( argumentativo ou expositivo ).
Esta competência envolve : - Estruturação sintática condizente com o padrão da modalidade escrita formal da Língua Portuguesa, de modo a garantir a clareza e a coerência necessárias ; - Distribuição do conteúdo do texto em parágrafos, de modo a garantir a sua organização temática ; - Distribuição das ideias em períodos que mantenham uma estruturação sintática condizente com a modalidade formal da língua ( evitando -se frases fragmentadas, sequência justaposta de ideias sem encaixamentos sintáticos e presença de oração subordinada sem oração principal que reflitam hábitos da oralidade e comprometam a estrutura lógico-gramatical do texto ) ; - Utilização de operadores discursivos que contribuam para a progressão temática do texto, estabelecendo relações lógicas entre as ideias apresentadas, tanto do ponto de vista intrafrasal, como interfrasal : preposições, conjunções, pronomes relativos, advérbios e locuções adverbiais de tempo e lugar ; - Utilização de procedimentos de referenciação lexical e pronominal que permitam a retomada de referentes textuais, como pronomes, advérbios e artigos ; além de elementos lexicais, como sinônimos, hiperônimos, hipônimos, expressões metafóricas e expressões resumitivas ) ; - Emprego adequado dos pronomes relativos “ cujo(a) ” e “ onde ”, responsáveis por encaixamentos sintáticos no período composto ; - Utilização adequada de sinais de pontuação que contribuam para a organização lógica da frase e do texto ; - Inteligibilidade relacionada ao atendimento das exigências de estruturação textual.
A expectativa, portanto, era de que o participante recorresse a procedimentos linguístico-discursivos para organizar seu texto, permitindo o encadeamento lógico entre suas partes de forma a garantir a progressão e a coerência textuais, com estruturação que não reproduzisse hábitos da oralidade e sim atendesse à norma-padrão da Língua Portuguesa. ( 2 ) Respeito às convenções ortográficas da norma-padrão da Língua Portuguesa.
Esta competência envolve o respeito às convenções estabelecidas pela legislação em vigor e consubstanciadas no Vocabulário Ortográfico da Língua Portuguesa, editado pela Academia Brasileira de Letras e avaliam -se : - domínio das regras de acentuação gráfica ; - domínio da grafia padrão das palavras ( com ausência de abreviaturas próprias da linguagem da internet ).
Assim, esperava -se que o participante : grafasse corretamente as palavras ; respeitasse as regras de acentuação gráfica ; empregasse maiúsculas em início de frase, em nomes próprios de pessoas, lugares ou instituições ; grafasse as siglas com letras maiúsculas ou apenas com a primeira letra maiúscula quando formarem uma palavra ; evitasse abreviações como p/, vc, tb, pq, \textbf{tá}, né, usadas muitas vezes em escrita informal e na internet ; obedecesse às regras de separação de sílabas no final da linha ; atendesse aos padrões de segmentação vocabular da norma-padrão. ( 3 ) Domínio dos diferentes aspectos morfossintáticos próprios da modalidade escrita formal da norma-padrão da Língua Portuguesa.
Esta competência envolve o domínio das convenções próprias da norma padrão da Língua Portuguesa relativas aos seguintes aspectos : concordância nominal ; concordância verbal ; regência nominal ; regência verbal ; flexão nominal ; flexão verbal ; correlação entre os tempos verbais ; colocação pronominal. ( 4 ) Assim, esperava -se que o participante : - flexionasse o verbo para estabelecer concordância de número com o sujeito da frase ; - flexionasse o artigo, o adjetivo e o pronome para concordar em número e em gênero com o substantivo a que \textbf{cada} um se refere ; - observasse a regência nominal e a verbal, utilizando a preposição adequada depois de um substantivo, um verbo ou um adjetivo ; - empregasse adequadamente o acento grave indicador da crase que expresse a contração entre uma preposição e um artigo, de acordo com as exigências de regência nominal ou verbal ; - obedecesse às regras de colocação pronominal ( próclise e ênclise ), distintas dos hábitos da oralidade ou da escrita informal ; - flexionasse adequadamente verbos, substantivos, adjetivos e pronomes, no que diz respeito à expressão das categorias gramaticais ; - flexionasse os verbos para expressar a correlação de modo e de tempo nas estruturas subordinadas.
Com base nesses critérios, alguns desvios ortográficos foram considerados como desvios de caráter morfossintático porque interferem na classificação dos vocábulos de base morfológica e funcional : - eliminação da marca de infinitivo ( -r- ) e substituição por acento agudo ou ausência total de marca do infinitivo, causando ambiguidade entre distintas classes de palavras ; - confusão entre as terminações “ ão ” e “ am ”, utilizadas para marcar diferentes tempos verbais ; - emprego indevido da forma verbal de"haver"\&lt ; “ há”\&gt ; no lugar da preposição ou do artigo \&lt ; “ a”\&gt ;, e vice-versa ; - emprego de hífen para separar pronome átono – tanto uso indevido quanto omissão ( ex : “ esperasse ”, em lugar de “ espera -se ” ; “ falar -mos ”, no lugar de “ falarmos ” ) ; - emprego de “ esta ” no lugar de “ \textbf{está} ”, provocando ambiguidade devido à confluência gráfica entre o pronome demonstrativo e a forma do verbo"\textbf{estar}", pela ausência da acentuação gráfica ; - emprego da conjunção aditiva"e"no lugar da forma verbal “ é ”, provocando ambiguidade devido à confluência gráfica por ausência da acentuação gráfica ; - emprego do advérbio “ mais ” no lugar da conjunção adversativa “ mas ”, reproduzindo confluência fônica da oralidade ; - registro da terceira pessoa do plural dos verbos “ ter ” e “ vir ” sem o acento circunflexo, causando desvio de concordância verbal ; - emprego inadequado do relativo “ onde ” deslocado do seu sentido locativo, reproduzindo tendência da oralidade ; - emprego do pronome relativo “ que ” sem a necessária preposição ( desvio de regência ) ; - ausência do acento grave obrigatório para marcar a crase. ( 5 ) Seleção vocabular adequada à modalidade escrita formal da Língua Portuguesa, exigida pela situação comunicativa.
Esta competência envolve : - precisão na seleção/utilização do vocabulário relacionado à temática solicitada pela questão ; - diversificação e adequação vocabular ; - coerência na combinação vocabular ; - ausência de marcas da oralidade, como termos de sentido muito genérico ( “ coisa ”, “ negócio ”, “ você ” ) e termos de registros mais informais ( como gírias, jargões, frases feitas, ditados populares, termos regionais ).
Assim, esperava -se que o participante : respeitasse a adequação vocabular, não empregando gírias ou expressões coloquiais ; apresentasse repertório vocabular diversificado e preciso ; evitasse repetição desnecessária de palavras ; utilizasse um vocabulário mais formal, como solicitado por um texto dissertativo.
Essas quatro competências foram reunidas em três itens de avaliação, atendendo aos aspectos explicitados anteriormente : I – Aspectos Ortográficos ( 30 \% ) : domínio das convenções ortográficas : grafia de vogais e consoantes, uso de maiúsculas e minúsculas, emprego do hífen e da acentuação gráfica, segmentação de palavras ; II – Aspectos Textuais ( 20 \% ) : domínio dos procedimentos de estruturação textual : organização interna dos períodos, emprego de conectores para a articulação lógica entre os períodos e entre os parágrafos, emprego de marcas de referenciação lexical e pronominal ; utilização dos sinais de pontuação que contribuem para a organização lógica da frase ; III – Aspectos Morfossintáticos e Vocabulares ( 50 \% ) : domínio das regras de caráter morfossintático estabelecidas como modelares do ponto de vista da modalidade escrita formal da norma-padrão da Língua Portuguesa : concordância nominal e verbal, regência nominal e verbal, colocação pronominal, flexão nominal e verbal, correlação entre tempos e modos verbais, ausência de marcas de oralidade.
Outra competência combinada a esta na análise das respostas é a seleção vocabular adequada à modalidade escrita formal da Língua Portuguesa, tendo em vista a intersecção entre as duas do ponto de vista das exigências do registro formal da modalidade escrita da norma-padrão. b ) Análise das respostas dadas pelos estudantes A Questão 1 propiciou a oportunidade de maior desenvolvimento do tema solicitado, gerando consequências na estruturação textual das respostas, já que os textos foram mais expositivos, mais elaborados e fluentes, inclusive, um dos textos era uma tabela.
Já a Questão 2 teve um comportamento distinto, devido ao seu caráter mais opinativo, o texto foi mais curto e o vocabulário mais limitado.
Observaram -se, nas duas questões, alguns casos de comprometimento de autoria, com textos muito próximos a paráfrases dos textos motivadores e sem marcas de repertório sociocultural produtivo.
Apresentam -se, a seguir, algumas observações sobre o desempenho dos participantes : I – Aspectos Ortográficos : o desempenho dos participantes revelou uma diferença muito grande nos dois aspectos analisados, como já foi observado nos relatórios dos anos anteriores : baixo índice de desvios da grafia padrão e grande índice de desvios de acentuação.
Em vários casos, ocorre ausência quase completa de acentuação gráfica.
Vale ressaltar que muitos participantes apresentaram uma caligrafia quase ininteligível.
Ainda nesse aspecto, ocorreu a omissão inadequada de cedilha, de til e do traço do T em algumas respostas.
Houve, inclusive, respostas com ausência total de cedilha ou til.
Os casos mais sistemáticos de eliminação do acento indicador da sílaba tônica foram : palavras proparoxítonas ( “ politica ”, “ pratica ”, “ publica ”,"policia") ; palavras paroxítonas terminadas em ditongo crescente ( “ necessario ”,"historia","propria", “ violencia ” ) ; palavras oxítonas ( “ tambem ”,"populacao", “ pais ” ).
Por outro lado, destaca -se o uso indevido do acento gráfico em palavras como “ investímento ” e “ pátriarcado ”, por exemplo.
Quanto aos desvios de grafia, observam -se os seguintes exemplos : emprego inadequado de maiúscula : “ Leis ” ; ou de minúscula :"constituição"“ Ministério publico ” ; confusão entre"ss","sc","c"e"ç": “ nessessaria", “ agreção ”, “ conciencia ”, “ comecam ”, “ cocientização ”, “ pacos ” ; eliminação de grafemas : “ mostando ”,"deveriamo"; - desvios de segmentação : “ des de ”, “ des igualdade ”, “ pre conceito ” ; troca de grafema : “ órgam ” ( por “ órgão ” ), “ escravido ” ( por “ escravidão ” ).
II – Aspectos Textuais : a estruturação textual foi a competência que apresentou o maior número de desvios, sendo responsável, em alguns textos, pela falta de clareza e pela desorganização.
No conjunto dos textos analisados, observaram -se diferentes níveis de domínio dos recursos de estruturação e organização textual.
No entanto, os participantes podem ser agrupados em três faixas descritas a seguir : - os que possuem grande domínio dos recursos de estruturação textual, com produção de textos que atendem às exigências de clareza e coerência, com progressão temática bem direcionada e utilização produtiva de operadores argumentativos ; - os que apresentam um domínio regular dos recursos de estruturação textual, com desvios na organização do texto e ausência de paragrafação, além de uso insuficiente de operadores argumentativos e processos de encaixamento sintático ; - os que apresentam baixo domínio dos recursos de estruturação textual, com muitos desvios que retomam hábitos da oralidade e que comprometem a coerência do texto : ocorrência de sequência justaposta de ideias sem encaixamentos sintáticos ; redução de estruturas subordinadas, ao lado do aumento na frequência de estruturas coordenadas e absolutas ; redução no uso de conectores para expressar relações lógicas essenciais à construção do texto ; emprego equivocado de operadores que não estabelecem relações lógicas coerentes entre ideias do texto ; fragmentação frasal que compromete a estrutura lógico-gramatical.
Uma estrutura que vem se ampliando na escrita dos participantes, refletindo hábitos da oralidade, é a duplicação de um dos elementos da frase, ou seja, a retomada anafórica de um elemento topicalizado, como por exemplo :"Elas precisam saber o verdadeiro papel DELAS na sociedade “, A desigualdade ELA vem de uma cultura machista... ” ; “ Homem e mulher, ELES são iguais ”.
Também em relação ao emprego dos sinais de pontuação, observaram -se distintos tipos de desempenho : participantes que dominam a funcionalidade da pontuação para a construção de textos fluentes e coesos, por um lado, e um grande número de participantes que apresentam precariedade em sua utilização, chegando ao ponto de elaborar parágrafos sem qualquer marca interna de pontuação para separar os períodos.
São os seguintes os tipos de desvios encontrados : - Emprego de vírgula para separar o sujeito e o predicado : “ A mulher, sofre violência ” ( no lugar de “ A mulher sofre violência ” ) ;"O correto seria, ter um canal de denúncia ” ( no lugar de “ o correto seria ter um canal de denúncia ” ) ; - ocorrência de apenas uma das vírgulas para separar uma palavra, uma expressão ou uma oração encaixada : “ O homem, por sua vez não aceita que uma mulher seja sua chefa ” ( no lugar de “ O homem, por sua vez, não aceita que uma mulher seja sua chefa ” ) ; - Emprego de ponto no lugar de vírgula para separar orações de um mesmo período : “ A todo custo as mulheres são atacadas na política.
Sexualmente também. ” ( no lugar de “ A todo custo as mulheres são atacadas na política, sexualmente também. ” ) ; - Ausência de vírgula ou conjunção para separar elementos de uma enumeração :"... conquistas como voto equidade salarial direito de estudar e reconhecimento..."( no lugar de""... conquistas como voto, equidade salarial, direito de estudar e reconhecimento...") ; - Ausência de vírgula para separar oração adjetiva explicativa ou utilização inadequada para separar oração adjetiva restritiva : “ As mulheres, que são iguais aos homens devem ser igualmente protegidas por lei."( no lugar de “ As mulheres, que são iguais aos homens, devem ser igualmente protegidas por lei.") ; - emprego indevido do sinal de dois pontos : “ Mulheres são julgadas : pelas simples roupas que usam. ” ( no lugar de “ Mulheres são julgadas pelas simples roupas que usam. ” ) ; - Ausência de pronome no processo de referenciação : “ julga a mulher ou até agride ” ( no lugar de “ julga a mulher ou até a agride ” ” ) ; - Emprego de vírgula em lugar do ponto-final : “ O patriarcado estrutural é formado por elementos históricos e religiosos, há pouquíssimo tempo, as mulheres passaram a votar ”. ( no lugar de “ O patriarcado estrutural é formado por elementos históricos e religiosos.
Há pouquíssimo tempo, as mulheres passaram a votar. ” ).
III – Aspectos Morfossintáticos e Vocabulares : o domínio das convenções morfossintáticas apresentou resultados que evidenciam, claramente, alguns aspectos problemáticos no desempenho dos participantes.
Assim como na análise dos desvios textuais, pode -se observar uma distinção entre três grupos de participantes : o grupo que revela domínio das convenções morfossintáticas ; o grupo que apresenta um domínio regular, com desvios pontuais ; e o grupo que revela baixo domínio dessas convenções ; e o grupo que revela domínio precário dessas convenções.
Assim como na análise dos desvios textuais, pode -se observar uma distinção entre três grupos de participantes : o grupo que revela domínio das convenções morfossintáticas ; o grupo que apresenta um domínio regular, com desvios pontuais ; e o grupo que revela domínio precário dessas convenções.
Pode -se destacar a regência, tanto verbal quanto nominal, como o aspecto com maior número de desvios : a ) Regência nominal e verbal : Na maioria dos textos, destacam -se os desvios relacionados à noção de"regência", ou seja, à relação de subordinação que se estabelece entre as palavras por meio, ou não, de preposições, refletindo uma tendência na modalidade escrita da língua.
O desvio mais frequente foi a falta do sinal indicativo da crase, em contextos em que seria obrigatória, ou, ao contrário, um uso indevido em contextos em que o princípio não se aplica.
Os dois tipos de desvios revelam que o usuário não tem consciência de que, sob a forma do termo “ a ”, se pressupõe a presença de uma preposição “ a ”, exigida pela regência do termo anterior.
Exemplos : “ diferentes tipos de agressão a mulher ” ( no lugar de “ diferentes tipos de agressão à mulher ” ) ; “ Desenvolverem práticas econômicas e sociais as novas gerações"( no lugar de “ Desenvolverem práticas econômicas e sociais às novas gerações") ; “ Além de proteger à mulher... ” ( no lugar de “ Além de proteger a mulher... ” ).
Outro desvio relacionado à regência verbal e à nominal, encontrado frequentemente nas questões, é a ausência de preposição antes de pronome relativo, refletindo um processo generalizado na modalidade oral da língua, em situações de registro informal.
Exemplo :"outra forma é oferecer apoio psicológico quem sofreu qualquer tipo de abuso..."( no lugar de :"outra forma é oferecer apoio psicológico a quem sofreu qualquer tipo de abuso...") b ) Concordância nominal e verbal : Os desvios relativos à aplicação das regras de concordância verbal e nominal correspondem a fatos que têm sido amplamente estudados na modalidade oral da Língua Portuguesa.
Quanto à concordância de número, observou -se ausência de marca ( com sujeito anteposto ou posposto ), como por exemplo a ausência de acento circunflexo na forma plural do presente do indicativo do verbo “ ter ”, considerada como um desvio na concordância verbal e não na acentuação gráfica.
Exemplo :"as mulheres ainda tem como obrigação"( no lugar de"as mulheres têm como obrigação") ; “ as leis em vigor tem que ser cumpridas ” ( no lugar de “ as leis em vigor têm que ser cumpridas ” ).
Outro desvio na aplicação da regra da concordância de número foi o aparecimento da marca de plural em verbos ou adjetivos relacionados a núcleos substantivos no singular, devido à presença de modificadores no plural, evidenciando um processo de hipercorreção.
Esse processo tem se ampliado nas situações de oralidade, caracterizando uma tendência de variação da Língua Portuguesa.
Exemplos :"Mudança nas leis que aumentem a pena do agressor"( no lugar de"Mudança nas leis que aumente a pena do agressor") ;"o aumento dos feminicídios devem ser combatidos ” ( no lugar de"o aumento dos feminicídios deve ser combatidos ” ).
Quanto à concordância de gênero, vários casos foram observados, normalmente no âmbito de sintagmas nominais longos, em que o adjetivo \textbf{está} afastado do substantivo, conforme apontam as pesquisas sobre a modalidade oral.
Exemplo : “ E as leis sobre as mulheres que sofrem violência ser mais rigorosos ” ( no lugar de “ E as leis sobre as mulheres que sofrem violência ser mais rigorosas ” ) ; Registro da forma verbal “ é ” sem acento, provocando ambiguidade na leitura pela confluência com a conjunção aditiva"e".
Exemplo : “ A desvantagem entre salários, poderes e entre a violência física e moral e crescente por demais... ” ( no lugar de “ A desvantagem entre salários, poderes e entre a violência física e moral é crescente por demais ”. c ) Emprego do pronome relativo “ onde ” como relativo universal, típica marca da oralidade, em que esse pronome não é mais utilizado apenas em sua função locativa"lugar onde se realiza a ação".
Exemplo :"Existe a desigualdade de gênero, onde homem ganha mais do que as mulheres"; “ há falta de informação e conhecimento onde ” ; “ O patriarcado \textbf{está} estruturado em nossa sociedade há séculos, onde o sexo masculino... ” d ) Falta de artigo definido antes de substantivo, comprometendo o paralelismo em enumerações : Exemplo :"seja em oportunidade de emprego, salário"( no lugar de"seja em oportunidade de emprego, de salário"). e ) Desvios relacionados à flexão verbal : Exemplos : emprego indevido do tempo ou do modo verbal : “ promovido ” por “ promovendo ” ; eliminação de marca de infinitivo : “ colocar ” por “ coloca ”, “ ta ” por “ \textbf{estar} ”, “ \textbf{está} ” por “ \textbf{estar} ”, “ esta ” por “ \textbf{estar} ” ; “ \textbf{estás} ” por “ estas ” ; utilização de construção de “ \textbf{estar} ” + gerúndio como marca de oralidade : “ a sociedade \textbf{cada} vez mais \textbf{tá} evoluindo ” ; segmentação indevida da desinência de número e pessoa de verbos : “ denunciar -mos ”.
Assim como nas demais competências avaliadas, há diferenças entre os participantes em relação à seleção e ao emprego do vocabulário.
Foram encontrados textos com alto padrão de domínio vocabular, chegando até ao emprego de vocabulário técnico das áreas do saber abordadas nas questões.
Por outro lado, na maioria dos textos, foram observados vários tipos de inadequação na seleção e na utilização do vocabulário, sendo que o principal deles foi a excessiva repetição de \textbf{certas} palavras, revelando limitação de repertório vocabular.
Além disso, observou -se a ocorrência de expressões da oralidade, revelando falta de domínio de vocabulário mais abstrato e de maior complexidade, essencial ao desenvolvimento do texto de base dissertativa.
Exemplos : seleção vocabular indevida : “ Em pleno século 21 é de emplena importância... ” ; emprego inadequado do verbo “ ter ” no lugar de “ haver ” : “ Durante a história, teve muito machismo ” ; imprecisão vocabular : “ Mas a cabeça, o corpo, há mesma igualdade para desenvolvermos igualmente as coisas ”, “ Em alguns anos atrás, as coisas eram mais complicadas nas questões de desigualdade no trabalho ”.
A imprecisão vocabular, como o uso da palavra “ coisa ” em vários textos, demonstra o domínio precário da norma-padrão da Língua Portuguesa, com presença constante de traços de oralidade. 7.3.1.8 Considerações finais sobre a correção de língua portuguesa das respostas às questões discursivas do componente de formação geral O conjunto de textos analisados no Enade 2022 permite traçar um perfil bem detalhado do desempenho linguístico do estudante brasileiro de graduação.
Conforme ficou comprovado nos comentários dos exemplos de respostas apresentados neste relatório, uma parcela significativa dos participantes revela lacunas importantes em sua formação, principalmente no que diz respeito à estruturação e à organização de textos.
Como esta é uma competência essencial à produção escrita, que deve ser construída ao longo dos anos de escolaridade ( do ensino fundamental à educação ), sua recuperação é trabalhosa e depende de um esforço intenso de exposição a textos e de reflexão sobre os requisitos que garantem a clareza e o encadeamento das ideias. competência essencial à produção escrita, que deve ser construída ao longo dos anos de escolaridade ( do ensino fundamental à educação ), sua recuperação é trabalhosa e depende de um esforço intenso de exposição a textos e de reflexão sobre os requisitos que garantem a clareza e o encadeamento das ideias.
A competência relacionada ao domínio do vocabulário, que deve ser diversificado e apropriado ao padrão formal da língua, também apresentou um desempenho insatisfatório pela maioria dos participantes, que se limitam ao senso comum, com vocabulário genérico e vago, o que é inadequado para atender ao perfil de um graduando.
As outras duas competências avaliadas ( convenções ortográficas, convenções morfossintáticas ) apresentam desvios pontuais e são de mais fácil recuperação, com exceção dos desvios morfossintáticos que já \textbf{estão} incorporados na modalidade oral informal da Língua Portuguesa, configurando processos de variação linguística e, talvez, antecipando futuras mudanças no padrão da modalidade escrita.
Com base na avaliação do desempenho linguístico aqui relatada, conclui -se que é preciso aperfeiçoar a formação superior, oportunizando aos estudantes mais acesso às ferramentas necessárias para que sejam bem-sucedidos nas diferentes situações de sua vida cidadã – profissional, cultural e política.
Essa é uma contribuição efetiva que os resultados do Enade podem oferecer às IES do país.
Para finalizar, vale observar que as duas questões discursivas atenderam às exigências de avaliação, tanto do domínio discursivo dos estudantes, quanto do seu conhecimento sobre temas relacionados ao conhecimento acadêmico, o que pode ser extremamente útil para que os cursos superiores possam contribuir na formação de seus graduandos. 7.3.2 COMPONENTE DE CONHECIMENTO ESPECÍFICO A Tabela 7.17 apresenta dados sobre questões discursivas do Componente Específico.
A partir dos dados registrados na Tabela 17.7, observa -se que a média foi mais baixa do que para as questões discursivas do Componente de Formação Geral.
Enquanto, no Componente de Formação Geral, a média para estudantes de Teologia de todo o Brasil foi 34,7 ; na parte de Conhecimento Específico, a média foi 34,7.
A maior média deste componente foi obtida pelos estudantes da região Norte ( 36,6 ), e a menor, pelos da região Nordeste ( 29,2 ).
Quanto à variabilidade das notas, o desvio-padrão de todo o Brasil foi 21,5.
O maior desvio-padrão foi encontrado na região Norte ( 23,0 ), e o menor, na região Centro-Oeste ( 22,2 ).
No Gráfico 7.10, apresenta -se a distribuição das notas nas questões discursivas no Componente Específico.
À luz dos dados do Gráfico 7.10, verifica -se que a moda da distribuição das notas nas questões discursivas no Componente Específico ocorre no intervalo ( 40 ; 50 ], com mais de 15 \% do total de participantes.
A análise de \textbf{cada} uma destas questões será feita a seguir. 7.3.2.1 ANÁLISE DA QUESTÃO DISCURSIVA 3 DO COMPONENTE DE CONHECIMENTO ESPECÍFICO Os resultados da Questão Discursiva 3 do Componente Específico são descritos na Tabela 7.18 e comentados na sequência.
Como mostrado na Tabela 7.18, na Questão 3, a média dos estudantes de todo o Brasil foi 38,6.
A menor média na questão, em evidência, foi obtida pelos estudantes da região Nordeste ( 32,5 ), enquanto a maior média foi obtida na região Norte ( 42,1 ).
Quanto à variabilidade das notas, o desvio-padrão de todo o Brasil foi 30,3.
O maior desvio-padrão foi obtido na região Norte ( 33,1 ), enquanto o menor foi obtido na região Nordeste ( 29,9 ).
A nota máxima, 100,0 pontos, foi alcançada por, pelo menos, um estudante em \textbf{cada} região.
A mediana do Brasil foi 40,0, para as regiões Sudeste e Sul ; nas demais regiões variou entre 33,3 e 46,7.
A nota mínima zero foi a mesma em todas as regiões do Brasil.
Destaca -se, com base nos dados do Gráfico 7.11, que mais de 23 \% dos estudantes deixaram a resposta à questão em branco.
A moda dessa distribuição, com exceção do intervalo [ 0 ; 10 ] ocorre no intervalo ( 50 ; 60 ], com mais de 15 \% do total de participantes. 7.3.2.2 COMENTÁRIOS SOBRE AS RESPOSTAS À QUESTÃO DISCURSIVA 3 A questão 3 abordou um tema relevante para a Teeologia, pois procurou explorar os conhecimentos relacionados ao marketing e às estratégias de atuação e comunicação.
Esperava -se, na estratégia de marketing, que o estudante demonstrasse conhecimentos sobre atuação para apresentação de novos produtos no mercado, como : identificar os produtos mais solicitados, fazer pesquisa junto aos clientes sobre produtos concorrentes que a empresa não tem similar e identificar inovações dos concorrentes que pudessem ser implementadas na empresa.
Na estratégia de comunicação desses produtos, que o estudante pudesse sugerir merchandising para chamar a atenção do consumidor ; promoção do produto no ponto de venda com degustação ou divulgação em redes e mídias sociais do novo produto.
O espaço para resposta era suficiente, o que chamou atenção foi que muitos estudantes não compreenderam o que o item “ a ” pedia e a resposta não contemplou uma pesquisa sobre a concorrência, em sua maioria.
Algumas provas apresentaram boas respostas, porém sem uma argumentação para dar sustentação.
Alguns estudantes conseguiram abordar os itens esperados para contemplar o recomendado para a questão, outros demonstraram uma \textbf{certa} uniformidade com o padrão de resposta.
No item “ a ”, percebeu -se que foram observadas, principalmente, as duas primeiras sugestões de resposta, nas quais se apresentou a possibilidade de uma pesquisa de mercado e divulgação em redes sociais, entretanto, as tendências de mercado e inovações foram pouco abordadas nas respostas.
Já no item “ b ”, todas as três alternativas de respostas foram exploradas.
Na questão “ a ”, nem todos os itens presentes no padrão de resposta foram encontrados.
O destaque foi para a aplicação da pesquisa.
Os alunos não se atentaram para a tendência de novos produtos e inovação dos concorrentes.
Nas respostas da questão “ b ”, o padrão estabelecido nas respostas foi o de merchandising, experimentação, promoção do produto in loco ou, ainda, divulgação nas redes e mídias sociais.
Na questão “ a ”, faltou explorar a atuação da concorrência sobre a estratégia junto aos clientes e o lançamento de novos produtos ( similares ) e, na questão “ b ”, poderiam ter explorado mais a questão do merchandising com cartaz ou display para divulgação do novo produto.
As respostas encontradas com maior frequência foram, na questão “ a ” : efetuar uma pesquisa de mercado ; na questão “ b ” : fazer a apresentação do produto através de degustação e divulgação em redes sociais.
Faltou explorar as tendências de mercado e inovações.
Os estudantes identificaram a necessidade da pesquisa junto aos clientes e a importância da demonstração do produto com degustação.
As respostas fracas foram aquelas que utilizaram uma palavra para sugerir a estratégia, sem explicar sua metodologia de aplicação.
Uma parte dos alunos demonstrou que domina o assunto, não apresentou nenhum tipo de desrespeito aos direitos humanos nos seus textos, foram isentos de opiniões pessoais e tentaram mostrar uma estratégia de atuação e comunicação.
As respostas \textbf{estão} dentro do padrão estipulado e do tipo de escrita e argumentação esperada por um aluno de graduação.
Boa parte dos alunos foi capaz de produzir textos dissertativos e condizentes com o enunciado da questão, associado aos conhecimentos adquiridos dentro de sala de aula.
Ficou evidente que os estudantes \textbf{estão} familiarizados com algumas estratégias para o lançamento de um novo produto, como a pesquisa de mercado, a utilização de redes sociais e a promoção de vendas através da utilização do promotor de venda apresentando o produto com degustação ou brindes.
Verificou -se que muitos alunos não discorreram sobre um dos pontos cobrados na avaliação : identificação de tendências de consumo sobre novos produtos na internet e redes sociais, como também inovações dos concorrentes de marcas possíveis de imitação.
Também muitos estudantes, apesar de colocarem palavras-chave, deixaram de explicar a estratégia para conseguir emplacar o novo produto. 7.3.2.3 ANÁLISE DA QUESTÃO DISCURSIVA 4 DO COMPONENTE DE CONHECIMENTO ESPECÍFICO Na Tabela 7.19, constam as informações relativas à Questão 4 do conjunto de Questões do Componente Específico.
Como se observa, a partir dos dados da Tabela 7.19, a média geral do Brasil foi 23,7, sendo a maior média registrada na região Norte ( 26,3 ), e a menor, na região Nordeste ( 21,0 ).
Quanto à variabilidade das notas, o desvio-padrão geral do Brasil foi 23,8 ; sendo o maior desvio-padrão registrado na região Norte ( 24,7 ), e o menor, na região Sul ( 23,5 ).
A nota máxima ( 100,0 ) foi atingida por, pelo menos, um concluinte nas regiões Centro-Oeste, Sudeste e Sul.
A mediana em todo o Brasil foi 20 ; o mesmo valor foi obtido nas regiões Centro-Oeste e Sul.
No Gráfico 7.12, \textbf{está} representada a distribuição de notas da Questão Discursiva 4, do Componente Específico.
A distribuição mostrada no Gráfico 7.12, excetuando -se resultados nulos, tem moda no intervalo ( 40 ; 50 ] com mais que 15 \% do total de participantes.
Chama a atenção que quase 40 \% dos participantes deixaram essa questão em branco. 7.3.2.4 COMENTÁRIOS SOBRE AS RESPOSTAS À QUESTÃO DISCURSIVA 4 O conhecimento avaliado pela Questão 4, para os estudantes da graduação, é básico e essencial para todas as atividades do curso.
A questão foi elaborada utilizando uma linguagem objetiva e clara, deixando poucas possibilidades de interpretações diferentes, podendo ser considerada de fácil construção, dado seu próprio enunciado e as informações apresentadas, possibilitando ao estudante condições de elaborar sua resposta.
A presença do gráfico remete ao caráter comparativo que uma resposta deve apresentar e, assim, estimula o aluno a trazer um contexto para encaixar os parâmetros das informações da questão.
A resposta do item “ a ” \textbf{está} praticamente toda contemplada no enunciado da questão, tornando -a fácil de ser respondida pelos estudantes, que assim o fizeram.
No item “ b ”, os estudantes responderam de forma similar ao sugerido no padrão.
Insistiram muito no conceito a respeito da confiabilidade, associando -o aos demais elementos, tais como preço e prazo de entrega.
Grande parte apenas citou esses fatores nas respostas ao item “ b ”, sem discorrer sobre seu significado no contexto da questão.
O tema faz parte da formação básica do curso e pertence às Diretrizes Curriculares Nacionais ( DCNs ) desde quando foi proposto e ofertado pela primeira vez.
Os estudantes, ao responderem à questão, inevitavelmente, utilizaram termos que já constavam no enunciado, especialmente na resposta ao item “ a ”.
O próprio enunciado já indicava o contexto da situação atípica de pandemia que nos assolou no período apresentado na questão.
Entretanto, muitos se ativeram mais em falar sobre a situação pandêmica e as restrições impostas do que sobre qualquer outra coisa.
Talvez por essa razão, grande parte foi rasa em suas respostas.
Já para as respostas ao item “ b ”, não foram utilizados os dados do enunciado, pois foi solicitada uma abordagem sobre o comportamento do consumidor na dimensão de seus fatores intrínsecos e extrínsecos, no que tange ao consumo presencial e virtual.
Apesar da questão trazer estes elementos, muitos apenas os citaram, sem o devido conceito contextualizado e sem o posicionamento do que influencia o comportamento do consumidor, que, no caso, são os mesmos fatores para ambos.
A linguagem adotada pelos estudantes, sobretudo nas respostas corretas, foi adequada para a reflexão e dissertação, conforme proposto pela questão.
Os estudantes, em boa parte das questões avaliadas, abordaram de forma correta a questão, indicando que a pandemia, apesar de ceifar vidas de maneira catastrófica, fez com que empresas entendessem que, mesmo diante de tais adversidades, de alguma forma, precisavam alcançar seus 7.3.2.5 ANÁLISE DA QUESTÃO DISCURSIVA 5 DO COMPONENTE DE CONHECIMENTO ESPECÍFICO Na Tabela 7.20, constam as informações relativas à Questão 5 do conjunto do Componente Específico.
Como revelam os dados da Tabela 7.20, o desempenho dos estudantes de todo o Brasil na Questão Discursiva 5 do Componente Específico trazendo nota média dos estudantes de todo o Brasil de consumidores e, por isso, o e-commerce potencialmente evoluiu.
Como a pergunta abordava os elementos que elevam a propensão no uso das ferramentas advindas do e-commerce, os estudantes responderam utilizando esses elementos como ponto de partida de suas respostas.
No item “ b ”, os estudantes, em sua maioria, citaram os elementos já expostos no referido item, no entanto, foram poucos que se aprofundaram nas descrições e relações destes elementos no comportamento do consumidor comparado às dimensões solicitadas, presencial e virtual.
Nas respostas aos itens “ a ” e “ b ”, as abordagens foram diretas e rasas, sem serem discursivas, especialmente no item “ b ”, em que lhes foi solicitado apenas citar.
Pode ser considerado que, de modo geral, ocorreu muita descrição da situação de pandemia, suas restrições, questões políticas, dentre outras coisas que notadamente ficaram em evidência na imprensa e nas mídias sociais e talvez por isso os alunos trouxeram este ponto para suas respostas.
As respostas dos estudantes demonstraram que a maioria conhece os elementos importantes sobre marketing, associados ao contexto da sociedade atual, entretanto, a objetividade excessiva das respostas evidenciou a falta de maior conhecimento de análise de dados, foco no tema e capacidade de organização das ideias de maneira coesa e coerente.
Esta última afirmação nos remete a uma possível deficiência de formação e cultura na área de dados.
Essa questão foi diretamente relacionada à formação do gestor comercial, essencial para aferir conhecimentos básicos do profissional. 7.3.2.5 ANÁLISE DA QUESTÃO DISCURSIVA 5 DO COMPONENTE DE CONHECIMENTO ESPECÍFICO Na Tabela 7.20, constam as informações relativas à Questão 5 do conjunto do Componente Específico. 41,6.
A maior média foi registrada na região Sudeste ( 43,9 ), enquanto a menor média foi registrada na região Nordeste ( 34,2 ).
Quanto à variabilidade das notas, o desvio-padrão dos estudantes do Brasil, como um todo, foi de 28,65.
Enquanto o maior desvio foi encontrado nas Regiões Sudeste e Norte ( 29,1 ), o menor foi encontrado na região Sul ( 27,9 ).
A mediana para o Brasil e para todas as regiões foi 46,7, excetuando -se Nordeste ( 33,3 ).
Para o conjunto de estudantes de Teologia do Brasil, a nota máxima foi 100,0, obtida por, pelo menos, um estudante em \textbf{cada} região.
A nota mínima do Brasil foi zero, o mesmo valor encontrado em todas as regiões.
No Gráfico 7.13, é representada a distribuição das notas da Questão Discursiva 5 do Componente Específico.
A distribuição mostrada no Gráfico 7.13, excetuando -se resultados nulos, tem moda no intervalo ( 60 ; 70 ], com mais de 20 \% do total.
Nessa questão, pouco mais de 25 \% a deixaram sem resposta. 7.3.2.6 COMENTÁRIOS SOBRE AS RESPOSTAS À QUESTÃO DISCURSIVA 5 O texto de Celidoni, T. \textbf{está} bem elaborado, claro, com linguagem de fácil compreensão ao estudante.
É possível que o aluno elabore uma boa resposta sobre o assunto, mesmo sem conhecimento prévio ; o embasamento teórico dado pelo texto norteia muito bem para que o discente discorra sobre o assunto.
É preciso salientar que o direcionamento final dado, ao solicitar que o candidato redija um texto, enfatiza que práticas publicitárias devem ser levadas em consideração e em nenhum momento anteriormente tal assunto foi levantado como critério de pontuação na construção textual.
A orientação aos corretores foi para pontuar com maior conceito candidatos que abordem as práticas publicitárias, como citado no comando da questão.
Profundidade, domínio do assunto e clareza na resposta devem ser critérios para padrão de resposta com pontuação máxima.
A instrução foi para considerar o real conhecimento que o candidato tem sobre o assunto principal, Lei Geral de Proteção de Dados ( LGPD ) – como a resposta foi desenvolvida sobre esse tema em específico, até onde vai a compreensão do aluno, para além do texto presente na questão, se as respostas \textbf{estão} se baseando apenas no texto de Celidoni, T., ou se o candidato já traz uma bagagem anterior sobre o assunto.
No mais, os critérios para o padrão de resposta contemplam o objetivo para avaliação de respostas satisfatórias, levantando o conhecimento do aluno sobre o tema.
Indica -se que seja considerada a sugestão dada sobre a “ apreciação geral da questão ”, texto acima, uma vez que as respostas obtidas pouco trouxeram o tema da prática publicitária no teor das respostas.
Seria interessante retirar as práticas de publicidade como direcionamento obrigatório de pontuação da questão, ou critério de pontuação máxima da mesma.
No mais, os critérios para padrão de resposta contemplam o intuito para avaliação de respostas satisfatórias.
Com frequência, as respostas obtidas não possuem muita profundidade, trazendo mais informações sobre a confidencialidade dos dados que a lei assegura.
As respostas, em sua maioria, vão nesse sentido antes e depois da implantação da LGPD, garantia da confidencialidade de dados dos compradores.
Algumas respostas indicam, ainda, como a tecnologia e as compras pela internet foram transformadas com a chegada dessa legislação, nada com muita profundidade ou bem fundamentado. 7.3.2.7 CONSIDERAÇÕES FINAIS SOBRE AS QUESTÕES DISCURSIVAS DO COMPONENTE DE CONHECIMENTO ESPECÍFICO Em primeiro lugar, é importante destacar que a avaliação do desempenho dos alunos em um exame como o Enade deve levar em conta diversos fatores, como o nível de dificuldade da prova, o perfil dos alunos que a habilitam e as condições em que a prova é aplicada.
Além disso, é preciso \textbf{lembrar} que o desempenho mediano não é necessariamente negativo.
Em uma distribuição normal, por exemplo, a maioria dos resultados se concentra em \textbf{torno} da média, o que significa que é esperado que um grande número de estudantes obtenha um desempenho mediano.
No entanto, se um aluno obteve um desempenho mediano, isso não significa que ele não possa melhorar.
É importante que o aluno faça uma reflexão sobre seus pontos fortes e fracos na prova e busque formas de aprimorar seu conhecimento e suas habilidades em áreas que ainda precisam de desenvolvimento.
É importante ampliar as discussões sobre como os estudantes de Teologia no Brasil \textbf{estão} recebendo sua formação.
As questões discursivas não permitiram diagnosticar se as informações recebidas nos cursos de graduação \textbf{estão} sendo prioritariamente técnicas, como em cursos técnicos, ou então se a formação dá ênfase a bases científicas para um maior aprofundamento e resoluções de problemas mais complexos.
Essa profissão é de cunho mais específico e analítico dentro das Ciências Após a correção de inúmeras respostas sobre a questão, é possível, seguramente, revelar : há uma falta considerável de conhecimento, em que, muitas vezes, há fuga daquilo que foi perguntado ; uma tentativa de respostas através do texto apresentado com preâmbulo da questão ; falta de interpretação de texto, já que houve um número elevado de respostas que não apresentaram nexo com o pedido da questão.
Ainda sobre a formação na área, podemos perceber que a bagagem individual, ou seja, o conhecimento para além do material trazido, é insuficiente ; a visão de resposta é limitada para apenas aquilo que foi \textbf{perguntado}, não conseguindo desenvolver uma resposta robusta e sólida.
A formação encontra -se prejudicada e teorizada, o que determina que o aluno não possui autonomia intelectual ampla, com visão estreita e reduzida para apenas o que foi solicitado.
Os níveis de respostas encontrados foram aquém do desejado.
Como dito anteriormente, foi preciso que os corretores focassem no conhecimento geral sobre o assunto abordado, pois a profundidade e a fundamentação na questão raramente aconteceram. 7.3.2.7 CONSIDERAÇÕES FINAIS SOBRE AS QUESTÕES DISCURSIVAS DO COMPONENTE DE CONHECIMENTO ESPECÍFICO Em primeiro lugar, é importante destacar que a avaliação do desempenho dos alunos em um exame como o Enade deve levar em conta diversos fatores, como o nível de dificuldade da prova, o perfil dos alunos que a habilitam e as condições em que a prova é aplicada.
Além disso, é preciso \textbf{lembrar} que o desempenho mediano não é necessariamente negativo.
Em uma distribuição normal, por exemplo, a maioria dos resultados se concentra em \textbf{torno} da média, o que significa que é esperado que um grande número de estudantes obtenha um desempenho mediano.
No entanto, se um aluno obteve um desempenho mediano, isso não significa que ele não possa melhorar.
É importante que o aluno faça uma reflexão sobre seus pontos fortes e fracos na prova e busque formas de aprimorar seu conhecimento e suas habilidades em áreas que ainda precisam de desenvolvimento.
É importante ampliar as discussões sobre como os estudantes de Teologia no Brasil \textbf{estão} recebendo sua formação.
As questões discursivas não permitiram diagnosticar se as informações recebidas nos cursos de graduação \textbf{estão} sendo prioritariamente técnicas, como em cursos técnicos, ou então se a formação dá ênfase a bases científicas para um maior aprofundamento e resoluções de problemas mais complexos.
Essa profissão é de cunho mais específico e analítico dentro das Ciências Sociais Aplicadas e seus profissionais devem receber informações científicas condizentes com sua importância.
O desempenho obtido nos remete a essa reflexão.
Por fim, é fundamental que o desempenho dos alunos no Enade seja utilizado como uma ferramenta para melhorar a qualidade do ensino.
Os resultados da prova devem ser analisados pelas instituições de ensino para identificar pontos fortes e falhas do currículo e implementar mudanças que possam contribuir para um melhor desempenho dos alunos no futuro.

% matched lemmas: cada, certo, estar, lembrar, ordem, perguntar, torno
\end{textsample}
